% 08_glossary.tex
\section{Glossary}\label{sec:glossary}

Plain-language definitions of the technical terms used in this guide, in alphabetical order.

\begin{description}[style=nextline, leftmargin=1.2em, itemsep=3pt, font=\normalfont\bfseries]

\item[Abadie--Imbens standard error]
A standard error for matching estimators that allows for the matching itself, including controls used more than once.

\item[AIPW]
See \emph{\hyperref[gl:aipw]{Augmented inverse-probability weighting (AIPW)}}.

\item[Area under the ROC curve (AUC)]\phantomsection\label{gl:auc}
A measure of discrimination: the probability that a model gives a higher predicted risk to a randomly chosen patient with the outcome than to one without it. 0.5 is no better than chance; 1 is perfect.

\item[Arm]
One of the groups compared in a trial or emulated trial, such as the treated arm and the untreated (control) arm.

\item[ATE]
See \emph{\hyperref[gl:ate]{Average treatment effect (ATE)}}.

\item[ATT]
See \emph{\hyperref[gl:att]{Average treatment effect in the treated (ATT)}}.

\item[AUC]
See \emph{\hyperref[gl:auc]{Area under the ROC curve (AUC)}}.

\item[Augmented inverse-probability weighting (AIPW)]\phantomsection\label{gl:aipw}
A doubly robust estimator that combines an outcome model with propensity-score weights. It corrects the outcome model's predictions using weighted residuals, so the estimate is correct if either model is correct.

\item[Average treatment effect (ATE)]\phantomsection\label{gl:ate}
The average effect if everyone in the population received the treatment, compared with no one receiving it.

\item[Average treatment effect in the treated (ATT)]\phantomsection\label{gl:att}
The average effect of treatment among the patients who actually received it: how their outcomes compare with what would have happened to them without treatment. Our default estimand.

\item[Balance]
How similar treated and untreated groups are on baseline characteristics, usually measured by the standardised mean difference.

\item[BART]
See \emph{\hyperref[gl:bart]{Bayesian additive regression trees (BART)}}.

\item[Bayesian additive regression trees (BART)]\phantomsection\label{gl:bart}
A Bayesian model that predicts the outcome as the sum of many small decision trees. Used to estimate treatment effects and their uncertainty (Section~\ref{sec:other_methods}).

\item[Benchmark bounds]
In sensemakr, the adjusted estimate if an unmeasured confounder were one, two or three times as strongly related to treatment and outcome as a named measured covariate.

\item[Best linear predictor (BLP) test]\phantomsection\label{gl:blp}
A check on heterogeneity: the doubly robust effect for each patient is regressed on the model's predicted individual effect. A slope above 0 means real variation in benefit has been detected; a slope near 1 means the individual predictions are well calibrated.

\item[Bias]
A systematic difference between an estimate and the true value that does not shrink as the sample grows.

\item[BLP]
See \emph{\hyperref[gl:blp]{Best linear predictor (BLP) test}}.

\item[Bootstrap]
Repeating the whole analysis on many samples drawn with replacement from the data, to estimate how much the result would vary from sample to sample. Used to build confidence intervals.

\item[Calibration]
Agreement between predicted probabilities and observed outcome rates. A well-calibrated model that predicts 30\% risk sees the outcome in about 30\% of such patients. Summarised by the calibration intercept (ideally 0) and slope (ideally 1).

\item[Caliper]
The largest difference allowed between a treated patient and their match, usually on the propensity score. Treated patients with no untreated patient within the caliper are left unmatched.

\item[CATE]
See \emph{\hyperref[gl:cate]{Conditional average treatment effect (CATE)}}.

\item[Causal forest]
A random forest adapted to estimate treatment effects rather than outcomes. Each tree splits patients into groups whose treatment effects differ.

\item[CausalForestDML]
The EconML class that fits a causal forest within the double machine learning framework.

\item[CI]
See \emph{\hyperref[gl:ci]{Confidence interval (CI)}}.

\item[CLAN]
See \emph{\hyperref[gl:gates]{Group average treatment effects (GATES) and classification analysis (CLAN)}}.

\item[Class weighting]
Giving more weight to the rarer outcome class when training a classifier. It improves accuracy-type metrics but distorts predicted probabilities, so it is avoided when effects are computed from probabilities.

\item[Clipping (of propensity scores)]
Replacing propensity scores below a small threshold (for example, 0.01) with that threshold, so that no patient receives an extreme weight.

\item[Clone-censor-weight]
A design for strategies defined by timing or duration: each patient is copied into every strategy they are compatible with at time zero, censored when they deviate from it, and the censoring is corrected by weighting. It avoids immortal-time bias.

\item[Cluster-robust confidence interval]
A confidence interval that allows observations within a cluster (for example, matched pairs sharing the same control) to be correlated.

\item[Coarsened exact matching]
Exact matching on coarsened (banded) versions of variables, which allows more variables to be matched while keeping most patients.

\item[Comparator]
The strategy that treatment is compared against, often ``no treatment''.

\item[Complier]
In an instrumental-variable analysis, a patient whose treatment depends on the instrument: they would be treated with one value of the instrument but not with another. Instrumental-variable estimates apply to compliers.

\item[Conditional average treatment effect (CATE)]\phantomsection\label{gl:cate}
The average effect for patients with a given set of characteristics. Estimated for each patient by the causal forest and the S-learner; also called the individual or heterogeneous treatment effect.

\item[Confidence interval (CI)]\phantomsection\label{gl:ci}
A range of values, calculated from the data, that would contain the true value in a stated proportion (usually 95\%) of repeated studies.

\item[Confounder]
A variable that affects both which treatment a patient receives and their outcome. Unadjusted confounding biases the estimated effect.

\item[Confounding by indication]
Confounding that arises because clinicians choose treatment based on the patient's prognosis or condition, so treated and untreated patients differ systematically.

\item[Consistency]
The assumption that ``treatment'' means the same thing for everyone given it, so that a patient's observed outcome equals their potential outcome under the treatment they received.

\item[Contour plot (sensitivity)]
In sensemakr, a plot of the adjusted estimate for every combination of a hypothetical confounder's partial $R^2$ with treatment and with outcome.

\item[Control (matched control)]
An untreated patient selected as the comparison for a treated patient.

\item[Control reuse]
Allowing one untreated patient to serve as the matched control for more than one treated patient (matching with replacement). It raises the matching rate but makes pairs less independent.

\item[Counterfactual (potential outcome)]
The outcome a patient would have had under a treatment they did not actually receive. It is never observed and must be estimated.

\item[Cross-fitting]
Splitting the data into folds and predicting each fold with a model trained on the other folds, so that no patient's prediction comes from a model that has seen them. It prevents overfitting from biasing effect estimates.

\item[Cross-validation]
Splitting the data into folds, training on some and testing on the rest in turn, to estimate how well a model performs on new data or to choose its settings.

\item[DAG]
See \emph{\hyperref[gl:dag]{Directed acyclic graph (DAG)}}.

\item[Decile]
One of ten equal-sized groups formed by ranking patients, for example by predicted risk.

\item[Delta method]
An approximation for the variance of a function of estimated quantities, such as the time at which predicted benefit reaches zero.

\item[Difference-in-differences]
Estimating the effect of a change by comparing the change in outcomes over time in places that made it with the change in places that did not.

\item[Directed acyclic graph (DAG)]\phantomsection\label{gl:dag}
A diagram of assumed causal relationships between variables, used to decide which variables to adjust for.

\item[Discrimination]
How well a model separates patients who have the outcome from those who do not, usually measured by the AUC.

\item[DML]
See \emph{\hyperref[gl:dml]{Double/debiased machine learning (DML)}}.

\item[Double/debiased machine learning (DML)]\phantomsection\label{gl:dml}
A framework that estimates causal effects using machine-learning nuisance models, with cross-fitting and residualisation so that errors in the nuisance models have only a small effect on the final estimate.

\item[Doubly robust]\phantomsection\label{gl:dr}
An estimator that remains consistent if either the outcome model or the treatment (propensity) model is correctly specified, not necessarily both.

\item[DR]
See \emph{\hyperref[gl:dr]{Doubly robust}}.

\item[DR-learner]
A meta-learner that estimates individual effects by regressing doubly robust pseudo-outcomes on patient characteristics.

\item[EconML]
A Python package from Microsoft Research implementing double machine learning, causal forests and meta-learners.

\item[Effect modifier]
A characteristic across which the size of the treatment effect varies (in the causal forest, the variables $X$).

\item[Eligibility criteria]
The conditions a patient must meet at time zero to be included in a trial or emulated trial.

\item[Entropy balancing]
A weighting method that chooses weights for untreated patients so that the means (and, optionally, variances) of the confounders exactly equal those of the treated patients.

\item[Estimand]
The precise quantity an analysis aims to estimate: which effect, in which population, on which scale (for example, the ATT as a difference in percentage points).

\item[E-value]
The minimum strength of association, on the risk-ratio scale, that an unmeasured confounder would need with both treatment and outcome to fully explain away an observed effect.

\item[Exact matching]
Pairing each treated patient with an untreated patient who has identical values on every matching variable.

\item[Exchangeability (no unmeasured confounding)]
The assumption that, after adjusting for measured variables, treated and untreated patients are comparable, so that the untreated represent what would have happened to the treated without treatment.

\item[Exclusion restriction]
The instrumental-variable assumption that the instrument affects the outcome only through the treatment.

\item[Extrapolation]
Making predictions for patients unlike any in the data used to fit the model, so that the estimate relies on modelling assumptions rather than data.

\item[First stage and second stage]
In double machine learning, the first stage fits the nuisance models (outcome and propensity); the second stage estimates the treatment effect from their residuals.

\item[Fold]
One of the equal parts into which data are split for cross-validation or cross-fitting.

\item[Follow-up]
The period from time zero over which outcomes are recorded.

\item[Full matching]
Matching that places every eligible patient in a subclass with at least one treated and one untreated patient, and uses weights, so that no eligible patient is discarded.

\item[GATES]
See \emph{\hyperref[gl:gates]{Group average treatment effects (GATES) and classification analysis (CLAN)}}.

\item[G-formula (g-computation)]
Estimating a treatment effect by modelling the outcome, predicting each patient's outcome under each treatment strategy, and averaging the predictions. The S-learner is a version of this for treatment at a single time.

\item[Gradient boosting]
A machine-learning method that builds many small decision trees in sequence, each correcting the errors of the previous ones (implemented in XGBoost and LightGBM).

\item[Group average treatment effects (GATES) and classification analysis (CLAN)]\phantomsection\label{gl:gates}
Group average treatment effects (GATES) compare average effects across groups defined by predicted benefit; classification analysis (CLAN) describes the characteristics of the most and least affected groups.

\item[Heterogeneity (of treatment effect)]
Variation in the size of the treatment effect between patients or subgroups.

\item[Holm adjustment]
A step-down method for controlling the chance of any false-positive finding when many hypotheses are tested.

\item[Honest tree]
A tree that uses one half of its sample to choose splits and the other half to estimate effects within the resulting groups, which reduces overfitting and supports valid confidence intervals.

\item[Hyperparameter tuning]
Choosing a model's settings (for example, tree depth or learning rate) by testing many combinations and keeping the one that performs best on held-out data.

\item[Immortal-time bias]
Bias that arises when a period during which the outcome cannot occur, such as the time before a treatment is given, is credited to the treated group.

\item[Influence function]
A description of how much each observation contributes to an estimate, used to calculate standard errors for doubly robust estimators.

\item[Instrumental variable]
A variable that affects which treatment is received but affects the outcome only through treatment (for example, a site's treatment rate). It can address unmeasured confounding under its own assumptions.

\item[Interference]
When one patient's treatment affects another patient's outcome. Most analyses assume there is none.

\item[Interrupted time series]
Estimating the effect of a change by comparing the trend in outcomes before and after it.

\item[Inverse probability weighting (IPW)]\phantomsection\label{gl:ipw}
Weighting each patient by the inverse of their probability of receiving the treatment they actually received, so that the weighted groups resemble each other.

\item[IPW]
See \emph{\hyperref[gl:ipw]{Inverse probability weighting (IPW)}}.

\item[Landmark analysis]
Restricting an analysis to patients still alive and in follow-up at a fixed time after time zero, as a crude check on immortal-time bias.

\item[Leave-one-out (of confounders)]
Re-estimating the effect with one measured confounder removed, to show how much the result depends on it.

\item[LightGBM]
A fast gradient-boosting library, used here for the first-stage models of the causal forest.

\item[Linear probability model]
A linear regression with a binary (0/1) outcome, whose coefficients are differences in probability. Used as a simple stand-in model for sensemakr.

\item[Logistic regression]
A regression model for a binary outcome that models the log-odds of the outcome as a linear combination of predictors.

\item[Log loss]
A measure of probabilistic prediction error that penalises confident wrong predictions. Used for tuning because it rewards both discrimination and calibration.

\item[Log-odds]
The logarithm of the odds, $\log\{p/(1-p)\}$. Effects on this scale are log odds ratios.

\item[Mahalanobis distance]
A measure of how far apart two patients are on several variables at once, allowing for the scale of each variable and the correlations between them.

\item[Marginal structural model]
A model of outcomes under different treatment strategies, fitted with inverse probability weights. Often used when treatments and confounders change over time.

\item[Matching rate]
The percentage of treated patients for whom a matched control was found.

\item[Matching variables]
The confounders on which treated and untreated patients must be identical to form a matched pair.

\item[McNemar's test]
A test for paired binary outcomes that compares the numbers of pairs in which only the treated, or only the control, patient had the outcome.

\item[Mediator]
A variable on the causal pathway between treatment and outcome. Adjusting for it removes part of the effect being estimated.

\item[Meta-learner]
A general recipe for estimating treatment effects using any prediction model; examples are the S-, T-, X- and DR-learners.

\item[Meta-model]
A model fitted to individual estimated benefits, used to describe which characteristics are associated with larger or smaller benefit.

\item[Missing as its own category]
Coding missing values of a categorical variable as a separate category, so that patients with missing data are kept in the analysis.

\item[Multiclass model]
A classification model that predicts the probability of each of more than two categories, such as several treatment strategies.

\item[Multiple imputation]
Filling in missing values several times by drawing from a model of the missing data, analysing each completed dataset and combining the results.

\item[Multiplicity]
The increased chance of false-positive findings when many statistical tests are carried out.

\item[Nearest-neighbour matching]
Matching each treated patient to the untreated patient closest to them on a distance measure, such as the propensity score or Mahalanobis distance.

\item[Negative control]
An outcome that the treatment cannot affect, or an exposure that cannot affect the outcome. An apparent effect on a negative control signals residual confounding or bias.

\item[Nested tuning]
Repeating hyperparameter tuning inside each training fold, so that the chosen settings never see the patients they are evaluated on.

\item[Non-exact matching]
Matching treated patients to untreated patients who are similar but not identical, using a distance measure, a propensity score or both (Section~\ref{sec:tte_nonexact}).

\item[Nuisance model]
A model needed to estimate the effect but not itself of interest, such as the outcome and propensity models in the first stage of double machine learning.

\item[Observational data]
Data collected during routine care or in registries, where treatment was not assigned at random.

\item[Odds ratio (OR)]\phantomsection\label{gl:or}
The ratio of the odds of the outcome in treated versus untreated patients. It is further from 1 than the risk ratio when the outcome is common.

\item[Omitted variable bias]
Bias in a regression estimate caused by leaving out a variable related to both treatment and outcome; the basis of sensemakr.

\item[Optimism]
The amount by which a model's performance on the data used to build it overstates its performance on new data.

\item[Optuna]
A Python library for hyperparameter tuning that searches efficiently and stops unpromising trials early.

\item[OR]
See \emph{\hyperref[gl:or]{Odds ratio (OR)}}.

\item[Out-of-bag estimate]
In a forest, a prediction for a patient made only by the trees that did not use that patient in their training sample.

\item[Out-of-fold prediction]
A prediction for a patient made by a model trained on other folds of the data (see cross-fitting).

\item[Overfitting]
When a model learns noise in the training data, so that it performs well on those data but poorly on new data.

\item[Overlap weights]
Propensity-score weights that give most weight to patients who could plausibly have received either treatment and little to those almost certain to receive one, which avoids extreme weights.

\item[Paired analysis]
An analysis that compares outcomes within matched pairs, such as McNemar's test, the paired $t$-test or the Wilcoxon signed-rank test.

\item[Parallel trends]
The difference-in-differences assumption that, without the change, outcomes in the compared groups would have followed the same path over time.

\item[Partial $R^2$]
The share of the remaining variance of a variable explained by another variable after accounting for the covariates; the scale used by sensemakr.

\item[Percentage points]
The unit of an absolute difference between two percentages (for example, 40\% versus 30\% is a difference of 10 percentage points).

\item[Placebo test]
An analysis in which the treatment or its timing is deliberately scrambled; any apparent effect points to bias in the method.

\item[Policy learning]
Using estimated treatment effects to derive and evaluate a rule for deciding whom to treat.

\item[Policy tree]
A shallow decision tree that assigns treatment based on a few patient characteristics, chosen to maximise estimated benefit.

\item[Positivity (overlap)]
The assumption that patients with each combination of characteristics have some chance of receiving each treatment, so comparable treated and untreated patients exist.

\item[Probability scale]
Expressing effects as differences in the probability of the outcome (in percentage points), as opposed to odds or log-odds.

\item[Propensity score]
A patient's probability of receiving a treatment given their characteristics.

\item[Propensity-score matching]
Matching treated and untreated patients on their propensity scores, rather than on each confounder separately.

\item[Protocol]
The written specification of a trial or emulated trial: eligibility, strategies, assignment, time zero, follow-up, outcomes, causal contrast and analysis.

\item[Pseudo-outcome]
A constructed outcome for each patient whose average equals the treatment effect, used as the target in a second-stage regression (for example, in the DR-learner).

\item[Quantitative bias analysis]
Estimating how much a specified source of bias, such as a named unmeasured confounder with assumed strength and prevalence, would change the result.

\item[Random forest]
A prediction method that averages many decision trees, each grown on a random sample of patients and features.

\item[Randomised controlled trial (RCT)]\phantomsection\label{gl:rct}
A study in which treatment is assigned by chance, so that treated and untreated groups are comparable on average in both measured and unmeasured characteristics.

\item[Random seed]
A number that fixes the sequence of random choices in an analysis, so that results can be reproduced.

\item[RCT]
See \emph{\hyperref[gl:rct]{Randomised controlled trial (RCT)}}.

\item[Regression adjustment]
Including covariates in a regression model of the outcome to remove remaining differences between groups.

\item[Regression discontinuity design]
Comparing patients just above and just below a threshold that decides treatment (such as a time limit), who are otherwise similar.

\item[Regularisation]
Penalising model complexity during fitting to reduce overfitting.

\item[Residual]
The difference between an observed value and a model's prediction of it.

\item[Residual confounding]
Confounding that remains after adjustment, for example because a variable was grouped into broad bands or measured imperfectly.

\item[Residualisation]
Removing the part of a variable that can be predicted from covariates, leaving only its residual.

\item[Risk ratio (RR)]\phantomsection\label{gl:rr}
The proportion with the outcome in treated patients divided by the proportion in untreated patients.

\item[Risk-set matching]
Matching each patient treated at a given time to a patient who is still untreated and still eligible at that time. Used when treatment strategies are defined by the timing of treatment.

\item[R-learner]
A meta-learner that estimates individual treatment effects from the residuals of the outcome and treatment models, which makes it robust to errors in those models.

\item[Robustness value (RV)]\phantomsection\label{gl:rv}
In sensemakr, the minimum partial $R^2$ an unmeasured confounder would need with both treatment and outcome to reduce the estimate to zero ($RV_\alpha$: to make it no longer statistically significant).

\item[ROC]
See \emph{\hyperref[gl:auc]{Area under the ROC curve (AUC)}}.

\item[RR]
See \emph{\hyperref[gl:rr]{Risk ratio (RR)}}.

\item[RV]
See \emph{\hyperref[gl:rv]{Robustness value (RV)}}.

\item[Selection by indication]
When the patients who undergo a test or receive a treatment are chosen because of their clinical features, so they differ systematically from those who do not.

\item[Sensemakr]
A sensitivity-analysis method that describes the strength of a hypothetical unmeasured confounder by its partial $R^2$ with treatment and outcome, benchmarked against measured covariates.

\item[Sensitivity analysis]
Repeating an analysis with one choice or assumption changed, to see whether the conclusions still hold.

\item[Sentinel value]
A value far outside the real range (for example, 99,999 minutes) used to code ``treatment not given'' in a continuous variable.

\item[SHAP]
See \emph{\hyperref[gl:shap]{SHapley Additive exPlanations (SHAP)}}.

\item[SHapley Additive exPlanations (SHAP)]\phantomsection\label{gl:shap}
A method that attributes a model's prediction for each patient to the contributions of individual features.

\item[S-learner (single learner)]
A meta-learner that fits one model to predict outcome from characteristics and treatment, then estimates effects by changing the treatment input.

\item[SMD]
See \emph{\hyperref[gl:smd]{Standardised mean difference (SMD)}}.

\item[SSNAP]
Sentinel Stroke National Audit Programme: the national audit of stroke care in England, Wales and Northern Ireland, and the source of the data analysed in the SAMueL-3 project.

\item[Standard deviation]
A measurement of the typical distance between individual numbers in a set and their average, showing how spread out the data is in the same units as the original numbers. To calculate it, find the variance, the average of the squared distances from the mean, and then take its square root.

\item[Standard error]
The estimated standard deviation of an estimate across repeated samples; used to build confidence intervals.

\item[Standardised mean difference (SMD)]\phantomsection\label{gl:smd}
The difference in means between two groups divided by their pooled standard deviation; absolute values below 0.1 usually indicate good balance.

\item[Stratum (plural strata)]
A subgroup of patients who share the same values of the stratifying (for example, matching) variables.

\item[Super Learner]
A method that combines several machine-learning models, weighted by their cross-validated performance. Often used with TMLE.

\item[Synthetic control]
A weighted combination of places (such as regions or hospitals) without a change, chosen to match a changed place's outcomes before the change, and used as its comparison afterwards.

\item[Targeted maximum likelihood estimation (TMLE)]\phantomsection\label{gl:tmle}
A doubly robust method that updates an initial outcome model in a targeted way to reduce bias in the effect estimate.

\item[TARGET statement]
Reporting guidance for observational studies that emulate a target trial.

\item[Target trial]
The hypothetical randomised trial that would answer the causal question, whose protocol an observational analysis is designed to emulate.

\item[Time zero]
The point at which eligibility is met, treatment strategies are assigned and follow-up begins. It must be the same for all compared groups.

\item[Time-varying confounding]
Confounding by factors that change during follow-up and may themselves be affected by earlier treatment.

\item[T-learner (two learner)]
A meta-learner that fits separate outcome models for treated and untreated patients and takes the difference in their predictions.

\item[TMLE]
See \emph{\hyperref[gl:tmle]{Targeted maximum likelihood estimation (TMLE)}}.

\item[Transportability]
Re-weighting an effect estimate so that it applies to a different target population, such as a trial population or another region.

\item[Treatment strategy]
A precisely defined plan of treatment, including what is given and when (for example, ``thrombolysis within 270 minutes'').

\item[Trimming]
Excluding patients with extreme propensity scores (for example, outside 0.05--0.95) so that estimates do not rely on extrapolation or extreme weights.

\item[TRIPOD]
Reporting guidance for studies that develop or validate prediction models.

\item[Unmeasured confounding]
Confounding by variables that are not recorded in the data, which no adjustment method can remove directly.

\item[Variance]
A measure how spread out a set of numbers is from their average, showing whether the values are tightly clustered near the middle or scattered far apart. To calculate it, find the average of the numbers, measure how far each number sits from that average, square those distances so negatives do not cancel out, and then take the average of those squared distances.

\item[Wilcoxon signed-rank test]
A non-parametric test comparing paired continuous measurements.

\item[XGBoost]
A widely used gradient-boosting library, used here for the S-learner outcome model.

\item[X-learner]
A meta-learner that imputes individual effects using models for each group, then combines them with weights based on the propensity score; useful when treatment groups differ greatly in size.

\end{description}
